%%% FIELD statement-rational-published
Under the premises of \(\texttt{thm:exact-once-output-sensitive-enumeration}\), suppose every entry of the rational input \(\mathcal O_G\) has binary encoding length at most \(L_{\mathrm{bits}}\). The fraction-free immutable-row implementation of \(\operatorname{Enum}_G\), including common-ancestry computation, has deterministic total bit-operation time and bit storage polynomial in \(p\), \(m\), \(L_{\mathrm{bits}}\), \(2^p\), \(N_{\mathrm{out}}\), and the total emitted encoding length. Canonically reduced determinant-ratio certificates decide equality of interned rows, and no superpolynomial intermediate or certificate growth is forced relative to those displayed parameters. Thus the guarantee is subset-exponential and output-sensitive, with the dependence on \(2^p\) explicit; it is not output-polynomial in the conventional input-plus-output sense.
%%% FIELD statement-rational-component-subset
Under the premises of \(\texttt{thm:component-aligned-exact-once-output-sensitive-enumeration}\), suppose every entry of the rational input \(\mathcal O_G\) has binary encoding length at most \(L_{\mathrm{bits}}\). A fraction-free implementation of \(\operatorname{Enum}_G^{\mathrm{CA}}\), including exact \(q\)-DAG construction and component-aligned common-ancestry computation, has a deterministic subset-exponential, output-sensitive total bit bound: its time and storage are polynomial in \(p\), \(m\), \(L_{\mathrm{bits}}\), \(2^p\), \(N_{\mathrm{out}}^{\mathrm{CA}}\), and the total emitted encoding length. This is not output-polynomial in the conventional input-plus-output sense because \(2^p\) remains explicit.
%%% FIELD proof-rational-published
Let \(\ell(z)\) denote the binary length of an integer \(z\), and represent each rational input coordinate as \(u/v\), where \(u,v\in\mathbb Z\), \(v>0\), \(\gcd(u,v)=1\), and \(\max\{\ell(u),\ell(v)\}\leq L_{\mathrm{bits}}\).  We first verify that every inverse and division used below is legitimate.  By \(\texttt{def:published-limvam-class}\), for a fixed view \(a\in M\) the disturbance coordinates are mutually independent and have strictly positive variances.  Hence, with
\[
 \Omega^{(a)}=\operatorname{Cov}_P(\varepsilon^{(a)}),
\]
the matrix \(\Omega^{(a)}\) is diagonal and positive definite.  By \(\texttt{ass:linear-structural-equations}\),
\[
 (I_p-B^{0,(a)})X^{(a)}=\varepsilon^{(a)}.
\]
The shared-order condition \(\texttt{ass:shared-topological-order}\) makes \(B^{0,(a)}\) strictly lower triangular after simultaneous permutation, so \(L_0^{(a)}=I_p-B^{0,(a)}\) is invertible.  Taking covariances gives
\[
 \Sigma^{aa}=(L_0^{(a)})^{-1}\Omega^{(a)}(L_0^{(a)})^{-\mathsf T}\succ0.
\tag{1}
\]
Consequently every principal matrix \(\Sigma^{aa}[A,A]\) in \(\texttt{def:immutable-original-row}\) is positive definite and nonsingular.  This also supplies the positivity and rank conditions required for the principal regression solves.

Fix one such solve with \(k=|A|\leq p\), and append its right-hand side to the \(k\times k\) coefficient matrix.  At most \((p+1)^2\) rational entries occur in this augmented matrix.  Let \(d\) be the product of their positive denominators.  Then
\[
 \ell(d)\leq (p+1)^2L_{\mathrm{bits}},
 \qquad
 \ell(dx)=O(p^2L_{\mathrm{bits}})
\tag{2}
\]
for every augmented entry \(x\).  Multiplication of the whole augmented system by the same nonzero scalar \(d\) leaves its solution unchanged.  After any fixed ordering of \(A\), all leading principal minors of the integer-scaled coefficient matrix are positive by (1), so no pivot is needed.

By the cited Bareiss result \(\texttt{lem:bareiss-minor-growth}\), the fraction-free recurrences make exact integer divisions and every transformed entry is a signed bordered minor of the scaled augmented matrix.  If \(K\) is an \(s\times s\) integer matrix with \(|K_{uv}|<2^h\), the Leibniz formula yields
\[
 |\det K|
 \leq \sum_{\sigma\in\mathfrak S_s}
       \prod_{u=1}^s |K_{u,\sigma(u)}|
 < s!\,2^{sh}.
\tag{3}
\]
Using \(s\leq p+1\) and (2), every Bareiss intermediate and every numerator or denominator supplied by the adjugate identity therefore has binary length
\[
 H_0=O(p^3L_{\mathrm{bits}}+p\log p).
\tag{4}
\]
The coefficient determinant is nonzero by (1), so the adjugate identity is applicable.  Dividing numerator and denominator by their greatest common divisor and choosing a positive denominator gives a unique rational certificate for every coordinate of \(\beta(A,r)\).  Thus two coordinates are equal if and only if their canonical numerator--denominator pairs agree.  Lexicographic order is also exact: for reduced fractions \(u/v\) and \(u'/v'\) with positive denominators, its scalar comparison is the sign comparison
\[
 uv'\mathrel{\lessgtr}u'v,
\]
whose operands have \(O(H_0)\) bits.  Equality and ordering of viewwise rows are obtained by finitely many such coordinate comparisons.  This proves the claimed deterministic row interning and verifies the certificate handle in \(\texttt{def:fraction-free-bit-handle}\).

Integer addition, multiplication, exact division, comparison, and greatest-common-divisor computation have bit cost polynomial in the operand lengths.  For completeness, Euclid's algorithm has only polynomially many remainder steps: after two consecutive nonterminal steps the positive second operand is less than half its value before those steps, so inputs of at most \(h\) bits require at most \(2h+1\) steps, each on operands of at most \(h\) bits.  Hence all row construction and interning operations have bit cost polynomial in \(p\) and \(L_{\mathrm{bits}}\).

By \(\texttt{thm:exact-once-output-sensitive-enumeration}\), the subsequent deterministic integer-trie computation emits exactly \(\Theta_G(\mathcal O_G)\), once per numerical tuple, and uses a number of ordered-exact-real operations and stored cells polynomial in \(p\), \(m\), \(2^p\), and \(N_{\mathrm{out}}\).  Fragment keys consist of at most \(p\) integer row identifiers, so their lengths are logarithmic in a quantity polynomial in the number of subset edges and cannot create rational-expression growth.  Decoding merely copies the canonical certificates already bounded in (4).  Published common ancestry is the empty relation by \(\texttt{prop:published-common-union-support-ancestry-empty}\), so that output requires \(O(1)\) further time and storage.

Multiplying the polynomial operation and cell counts by the polynomial bit cost and height bounds just proved, and adding the total length of the emitted certificates, gives deterministic total bit time and storage polynomial in
\[
 p,\quad m,\quad L_{\mathrm{bits}},\quad 2^p,
 \quad N_{\mathrm{out}},\quad
 \text{and the total emitted encoding length}.
\]
Every regression row is solved directly from an original principal covariance system, so (4) is not recursively compounded along a subset path.  Finally, the displayed parameter list retains \(2^p\).  Since \(2^p\) need not be polynomial in the rational input length plus the output length when, for example, \(N_{\mathrm{out}}=1\), the proved guarantee is subset-exponential and output-sensitive, not output-polynomial in the conventional input-plus-output sense.
%%% FIELD proof-rational-component-subset
We verify first the positivity and height conditions needed by the component-aligned preprocessing.  Fix a maintained mechanism and a view \(a\in M\).  By \(\texttt{ass:component-aligned-disturbances}\), distinct coordinates of \(\varepsilon^{(a)}\) are uncorrelated.  By \(\texttt{ass:disturbance-positive-definiteness}\),
\[
 \operatorname{Var}_P(\varepsilon_j^{(a)})
 =\Gamma_j[a,a]>0
 \qquad (j\in V).
\tag{5}
\]
Thus \(\Omega^{(a)}=\operatorname{Cov}_P(\varepsilon^{(a)})\) is diagonal and positive definite.  The structural equation and shared-order assumptions give, exactly as in (1),
\[
 \Sigma^{aa}=(I_p-B^{0,(a)})^{-1}
 \Omega^{(a)}(I_p-B^{0,(a)})^{-\mathsf T}\succ0.
\tag{6}
\]
Therefore every principal system used in an immutable row or projected residual state is nonsingular.

For \(A\subseteq V\), write
\[
 S_A^{(a)}=Sigma^{aa}[U_A,A]
 \{\Sigma^{aa}[A,A]\}^{-1},
\]
with the evident zero matrix when \(A=\varnothing\).  The denominator-clearing, Bareiss, determinant bound, and canonical reduction proved in \(\texttt{thm:rational-output-polynomial-bit-enumeration}\) apply verbatim to every coordinate of \(S_A^{(a)}\) and every immutable row \(\beta(A,r)\), because their coefficient matrices are the positive-definite principal matrices in (6).  In particular these coordinates have binary height polynomial in \(p\) and \(L_{\mathrm{bits}}\).

Expanding the covariance of the two residuals in \(\texttt{def:projected-residual-state}\) by bilinearity gives, on every observed diagonal block and oriented edge \((a,b)\),
\[
 \begin{aligned}
 C_A^{ab}
 ={}&\Sigma^{ab}[U_A,U_A]
 -S_A^{(a)}\Sigma^{ab}[A,U_A]
 -\Sigma^{ab}[U_A,A](S_A^{(b)})^{\mathsf T}\\
 &+S_A^{(a)}\Sigma^{ab}[A,A](S_A^{(b)})^{\mathsf T}.
 \end{aligned}
\tag{7}
\]
This is a direct formula in the original input blocks and two principal-system solutions; no state \(C_{A'}\) for a smaller subset is substituted into it.  Each coordinate of (7) is a sum of at most a polynomial number of products of polynomial-height rationals.  Repeated common-denominator multiplication can increase binary length by at most the number of summands times the largest operand length.  Hence every coordinate of every \(C_A^{ab}\) has binary height polynomial in \(p\) and \(L_{\mathrm{bits}}\), uniformly over the \(2^p\) subsets.

For the division in the root score, the within-view specialization of (7) is the Schur complement
\[
 C_A^{bb}
 =\Sigma^{bb}[U_A,U_A]
 -\Sigma^{bb}[U_A,A]
  \{\Sigma^{bb}[A,A]\}^{-1}
  \Sigma^{bb}[A,U_A].
\tag{8}
\]
The Schur complement of a principal block of the positive-definite matrix in (6) is positive definite: indeed, for nonzero \(y\in\mathbb R^{U_A}\), set
\[
 x_A=-\{\Sigma^{bb}[A,A]\}^{-1}
       \Sigma^{bb}[A,U_A]y,
 \qquad x_{U_A}=y.
\]
Then direct block multiplication yields
\[
 y^{\mathsf T}C_A^{bb}y
 =x^{\mathsf T}\Sigma^{bb}x>0.
\tag{9}
\]
In particular \(C_A^{bb}[r,r]>0\) for every \(r\in U_A\).  The division in \(\texttt{def:q-root-score}\) is therefore valid, and
\[
 q^{ab}(A,r,t)=0
 \quad\Longleftrightarrow\quad
 C_A^{ab}[r,t]C_A^{bb}[r,r]
 -C_A^{ab}[r,r]C_A^{bb}[r,t]=0.
\tag{10}
\]
Both sides of the integer cross multiplication in (10) have polynomial binary height by (7).  Thus every edge decision in \(\mathcal D_G\), both orientations included, is an exact polynomial-bit-cost decision.  Source-to-sink pruning then uses only the resulting finite Boolean graph.  The maintained-model support and positive-definite-completion validity of this graph are supplied by \(\texttt{thm:qpath-completion-exactness}\); the exact coefficient projection and exact-once output statement are supplied by \(\texttt{thm:component-aligned-exact-once-output-sensitive-enumeration}\).

The latter theorem bounds the number of exact-real operations by
\[
 O\!\left((m+|E|)p^3 2^p
 +p^3 2^pN_{\mathrm{out}}^{\mathrm{CA}}
 +mp^2N_{\mathrm{out}}^{\mathrm{CA}}\right)
\tag{11}
\]
and bounds storage by a polynomial in \(p\), \(m\), \(2^p\), and \(N_{\mathrm{out}}^{\mathrm{CA}}\).  Equations (5)--(10) show that every rational arithmetic operation or comparison in the direct Schur-state and \(q\)-DAG preprocessing has polynomial-height operands and polynomial bit cost.  The cited Bareiss lemma and the published rational theorem give the same conclusion for immutable-row certificates.  Fragment propagation uses only canonical integer identifiers.  Exact support tests on decoded rational rows, Boolean transitive closure, and intersection compute component-aligned common ancestry within the last term of (11), without increasing rational height.

Multiplying (11) and the corresponding storage count by the polynomial per-operation bit cost, and adding the total emitted encoding length, proves that total bit time and storage are polynomial in
\[
 p,\quad m,\quad L_{\mathrm{bits}},\quad 2^p,
 \quad N_{\mathrm{out}}^{\mathrm{CA}},\quad
 \text{and the total emitted encoding length}.
\]
Because every state is evaluated directly by (7), the height estimate is not recursively exponentiated through the subset lattice.  The factor \(2^p\) nevertheless remains explicit even when the numerical output set is small.  Thus this is a deterministic subset-exponential, output-sensitive total bit bound, not an output-polynomial bound in the conventional input-plus-output sense.
%%% FIELD proof-published-ancestry-empty
By \(\texttt{thm:published-all-orders-covariance-compatible}\), the published coefficient projection is nonempty and satisfies
\[
 \Theta_G(\mathcal O_G)=\{T_\pi:\pi\in\mathfrak S_V\}.
\tag{12}
\]
Fix \(\pi\in\mathfrak S_V\).  In \(\texttt{def:modified-cholesky}\), the only entries that can be nonzero in row \(\pi(k)\) have column labels in \(\{\pi(1),\ldots,\pi(k-1)\}\).  Hence every union-support edge \(u\to v\) of \(\mathcal U(T_\pi)\) points from an earlier vertex to a later vertex in \(\pi\).  Along every directed path the position in \(\pi\) therefore increases strictly.  It follows both that \(\mathcal U(T_\pi)\) is acyclic and that its strict transitive closure contains no diagonal pair.

Now fix distinct \(i,j\in V\).  Choose a permutation \(\pi_{j<i}\) in which \(j\) occurs before \(i\); such a permutation exists because \(V\) is finite and contains both labels.  Equation (12) gives \(T_{\pi_{j<i}}\in\Theta_G(\mathcal O_G)\).  Since every directed path in its union-support graph moves forward in \(\pi_{j<i}\), there is no path from the later vertex \(i\) to the earlier vertex \(j\).  Therefore
\[
 (i,j)\notin
 \operatorname{TC}^{+}(\mathcal U(T_{\pi_{j<i}})).
\tag{13}
\]
The definition \(\texttt{def:common-union-support-ancestry}\) intersects these strict transitive closures over all \(T\in\Theta_G(\mathcal O_G)\).  Equation (13) excludes every off-diagonal ordered pair from that intersection, while acyclicity excludes every diagonal pair.  Consequently
\[
 \mathcal A_G(\mathcal O_G)=\varnothing.
\]
The algorithm may therefore return the fixed empty relation without inspecting the tuple stream; writing this predeclared result takes \(O(1)\) time and \(O(1)\) space.
%%% FIELD tldr-update
The observed covariance law supports two exact but different conclusions. Under the published shared-order LiMVAM primitives, unrestricted cross-view disturbance dependence makes every order covariance-compatible and makes common union-support ancestry empty, so pairing contributes neither pruning nor ancestry. Under the component-aligned subclass, the independently defined mechanism fiber is exactly represented by source-to-sink paths of the live \(q\)-DAG. A two-mode immutable-row dynamic program emits each distinct coefficient tuple once; its pairing-informative component-aligned mode runs in \(O((m+|E|)p^3 2^p+p^3 2^pN_{\mathrm{out}}^{\mathrm{CA}}+mp^2N_{\mathrm{out}}^{\mathrm{CA}})\) ordered exact-real time and computes the nontrivial subclass ancestry intersection. For rational inputs, both modes have deterministic subset-exponential, output-sensitive total bit bounds polynomial in the listed parameters including explicit \(2^p\); these are not output-polynomial guarantees in the conventional input-plus-output sense.
%%% FIELD technical-limitation-update
The guarantees are total output-sensitive bounds over an explicit \(2^p\) state space, not polynomial delay, not polynomial time in \(p\), and not output-polynomial in the conventional input-plus-output sense. Ordered exact-real arithmetic is not a floating-point stability guarantee; for rational inputs the bit theorems instead control certificate height and total bit cost as polynomials in parameter lists that explicitly include \(2^p\). The \(q\)-DAG theorem requires maintained component-aligned inputs and does not provide polynomial preprocessing for arbitrary off-model nonchordal partial covariances. Gaussianity is used only to construct existential completion witnesses, not assumed of the generating law. The results enumerate exact population coefficient sets, not continuum-valued covariance-region inverses.
%%% FIELD honest-scope-update
The published shared-order class retains arbitrary cross-view disturbance covariance: pairing supplies neither order pruning nor common ancestry there, although the full-lattice algorithm still removes numerical duplicates without enumerating permutations. The pairing-informative result is exactly the component-aligned subclass with positive-definite coordinate covariance blocks, on an arbitrary observed view graph and at model-generated population moments. It proves \(q\)-path/fiber equivalence, exact-once tuple enumeration, nontrivial ancestry post-processing, ordered-exact-real complexity, and deterministic rational subset-exponential, output-sensitive total bit complexity with explicit \(2^p\) dependence. It does not claim unrestricted LiMVAM identification, conventional output-polynomial complexity, polynomial delay, off-model nonchordal completion, floating-point robustness, or finite enumeration of the parent's covariance-region inverse.
